\section{Ordinary Petz overlap under quantum channels}
\label{sec:petz_overlap}

\begin{definition}[Ordinary Petz overlap]
    \label{def:petz_overlap}
    \lean{Matrix.petzOverlap}
    \leanok
    For positive semidefinite matrices $\rho,\sigma$ on the same finite
    coordinate space and $\alpha\in[0,1]$, set
    \begin{align}
        Q_\alpha(\rho,\sigma)
            &=\operatorname{Re}\tr(\rho^\alpha\sigma^{1-\alpha}).
        \label{eq:petz_overlap_definition}
    \end{align}
    Powers are defined by spectral calculus, with $x^0=1$ also at $x=0$.
    In particular, the zeroth power of a singular matrix is the identity.
\end{definition}

\begin{lemma}[Endpoint values and invariance]
    \label{lem:petz_overlap_invariance}
    \lean{Matrix.petzOverlap_zero, Matrix.petzOverlap_one,
        Matrix.petzOverlap_conj_unitary, Matrix.petzOverlap_reindex,
        Matrix.petzOverlap_tensor_density}
    \leanok
    \uses{def:petz_overlap}
    For positive semidefinite $\rho,\sigma$, the endpoint values are
    $Q_0(\rho,\sigma)=\operatorname{Re}\tr\sigma$ and
    $Q_1(\rho,\sigma)=\operatorname{Re}\tr\rho$.
    Simultaneous unitary conjugation and changes of finite coordinates
    preserve $Q_\alpha$. If $\tau$ is any density matrix, then
    \begin{align}
        Q_\alpha(\rho\otimes\tau,\sigma\otimes\tau)
            &=Q_\alpha(\rho,\sigma).
        \label{eq:petz_overlap_ancilla}
    \end{align}
    Singular states and singular ancillas are permitted.
\end{lemma}
\begin{proof}
    \leanok
    Spectral calculus gives the endpoint identities and commutes with
    unitary conjugation and changes of coordinates. Cyclicity of the trace
    then proves the corresponding invariances. Powers factor across tensor
    products of positive semidefinite matrices. Since $\alpha$ and $1-\alpha$
    are nonnegative, the scalar identity
    $x^\alpha x^{1-\alpha}=x$ holds for $x\ge0$, including $x=0$.
    Therefore $\tau^\alpha\tau^{1-\alpha}=\tau$, and the tensor trace
    factors with multiplier $\tr\tau=1$.
\end{proof}

\begin{theorem}[Data processing for the ordinary Petz overlap]
    \label{thm:petz_overlap_data_processing}
    \lean{Matrix.petzOverlap_le_map, Matrix.petzOverlap_data_processing,
        Matrix.petzOverlap_le_partialTraceRight}
    \leanok
    \uses{def:petz_overlap}
    Let $T:\MN d\to\MN{d'}$ be completely positive and trace preserving,
    let $\rho,\sigma$ be density matrices, and let $\alpha\in[0,1]$. Then
    \begin{align}
        Q_\alpha(\rho,\sigma)
            &\le Q_\alpha(T(\rho),T(\sigma)).
        \label{eq:petz_overlap_data_processing}
    \end{align}
    This is \cite[Equation~(8.35)]{Wolf2012Quantum} on the stated exponent
    interval. More generally the same inequality holds for arbitrary positive
    semidefinite $\rho,\sigma$ when $d\ge1$, without trace normalization.
    In particular, discarding a nonzero right tensor factor increases the
    overlap; the retained coordinate space may have dimension zero.
    The reverse inequality for negative exponents is a separate assertion.
\end{theorem}
\begin{proof}
    \leanok
    \uses{thm:lieb_concavity_psd, lem:petz_overlap_invariance,
        thm:convex_functional_monotonicity, thm:convex_functional_partial_trace}
    Lieb's theorem with the auxiliary matrix equal to the identity makes
    $-Q_\alpha$ jointly convex on pairs of positive semidefinite matrices.
    The domain is convex and closed under simultaneous unitary conjugation,
    tensoring with a density matrix, and changes of coordinates.
    Lemma~\ref{lem:petz_overlap_invariance} proves invariance of the value.
    A completely positive map preserves positivity, so the image pair belongs
    to the output domain. Theorem~\ref{thm:convex_functional_monotonicity}
    applied to $-Q_\alpha$ proves the inequality. A trace-one input supplies
    $d\ge1$. The partial-trace assertion follows from
    Theorem~\ref{thm:convex_functional_partial_trace} on the same domain.
\end{proof}
